\documentclass{article}

% packages
\usepackage{amsmath, amsthm, thmtools, amsfonts, amssymb, luacode, catchfile, tikzducks, hyperref, ifthen}
\ifcsname c@kobocompile\endcsname
	\usepackage[a5paper, total={1072pt, 1448pt}, margin=10pt, includeheadfoot]{geometry} % set page margins
\else
	\usepackage[a4paper, margin=50pt, includeheadfoot]{geometry}
\fi
\usepackage[shortlabels]{enumitem}
\usepackage[skip=3pt, indent=0pt]{parskip}

% language
\usepackage[bidi=basic, layout=tabular, provide=*]{babel}
\ifcsname c@english\endcsname
	\babelprovide[main, import]{english}
\else
	\babelprovide[main, import]{hebrew}
	\babelprovide{rl}
\fi
%\babelfont{rm}{Libertinus Serif}
\babelfont{rm}[Renderer=Harfbuzz]{Libertinus Serif}
\babelfont{sf}{Libertinus Sans}
\babelfont{tt}{Libertinus Mono}

% style
\AddToHook{cmd/section/before}{\clearpage}	% Add line break before section
\linespread{1.3}
\setcounter{secnumdepth}{0}		% Remove default number tags from sections, this won't do well with theorems
\AtBeginDocument{\setlength{\belowdisplayskip}{3pt}}
\AtBeginDocument{\setlength{\abovedisplayskip}{3pt}}
\graphicspath{ {../images/} }

% operators
\DeclareMathOperator\cis{cis}
\DeclareMathOperator\Sp{Sp}
\DeclareMathOperator\tr{tr}
\DeclareMathOperator\im{Im}
\DeclareMathOperator\re{Re}
\DeclareMathOperator\diag{diag}
\DeclareMathOperator*\lowlim{\underline{lim}}
\DeclareMathOperator*\uplim{\overline{lim}}
\DeclareMathOperator\rng{rng}
\DeclareMathOperator\Sym{Sym}
\DeclareMathOperator\Arg{Arg}
\DeclareMathOperator\Log{Log}
\DeclareMathOperator\dom{dom}
\DeclareMathOperator\supp{Supp}
\DeclareMathOperator\var{Var}
\DeclareMathOperator\cov{Cov}

% commands
%\renewcommand\qedsymbol{\textbf{מש''ל}}
%\renewcommand\qedsymbol{\fbox{\emoji{lizard}}}
\newcommand{\Aa}[0]{\mathcal{A}}
\newcommand{\Bb}[0]{\mathcal{B}}
\newcommand{\CC}[0]{\mathbb{C}}
\newcommand{\Cc}[0]{\mathcal{C}}
\newcommand{\EE}[0]{\mathbb{E}}
\newcommand{\FF}[0]{\mathbb{F}}
\newcommand{\Ff}[0]{\mathcal{F}}
\newcommand{\Ii}[0]{\mathcal{I}}
\newcommand{\Gg}[0]{\mathcal{G}}
\newcommand{\Ll}[0]{\mathcal{L}}
\newcommand{\Mm}[0]{\mathcal{M}}
\newcommand{\NN}[0]{\mathbb{N}}
\newcommand{\Nn}[0]{\mathcal{N}}
\newcommand{\PP}[0]{\mathbb{P}}
\newcommand{\Pp}[0]{\mathcal{P}}
\newcommand{\QQ}[0]{\mathbb{Q}}
\newcommand{\RR}[0]{\mathbb{R}}
\newcommand{\Rr}[0]{\mathcal{R}}
\newcommand{\Ss}[0]{\mathcal{S}}
\newcommand{\TT}[0]{\mathbb{T}}
\newcommand{\Uu}[0]{\mathcal{U}}
\newcommand{\Vv}[0]{\mathcal{V}}
\newcommand{\Ww}[0]{\mathcal{W}}
\newcommand{\ZZ}[0]{\mathbb{Z}}
\newcommand{\acts}[0]{\circlearrowright}
\newcommand{\explain}[2] {
	\begin{flalign*}
		 && \text{#2} && \text{#1}
	\end{flalign*}
}
\newcommand{\maketitleprint}[0]{ \begin{center}
	%\begin{tikzpicture}[scale=3]
	%	\duck[graduate=gray!20!black, tassel=red!70!black]
	%\end{tikzpicture}	
	\includegraphics[width=6cm]{cover}
\end{center}
}

% theorem commands
\newtheoremstyle{c_remark}
	{}	% Space above
	{}	% Space below
	{}% Body font
	{}	% Indent amount
	{\bfseries}	% Theorem head font
	{}	% Punctuation after theorem head
	{.5em}	% Space after theorem head
	{\thmname{#1}\thmnumber{ #2}\thmnote{ \normalfont{\text{(#3)}}}}	% head content
\newtheoremstyle{c_definition}
	{3pt}	% Space above
	{3pt}	% Space below
	{}% Body font
	{}	% Indent amount
	{\bfseries}	% Theorem head font
	{}	% Punctuation after theorem head
	{.5em}	% Space after theorem head
	{\thmname{#1}\thmnumber{ #2}\thmnote{ \normalfont{\text{(#3)}}}}	% head content
\newtheoremstyle{c_plain}
	{3pt}	% Space above
	{3pt}	% Space below
	{\itshape}% Body font
	{}	% Indent amount
	{\bfseries}	% Theorem head font
	{}	% Punctuation after theorem head
	{.5em}	% Space after theorem head
	{\thmname{#1}\thmnumber{ #2}\thmnote{ \text{(#3)}}}	% head content

\ifcsname c@english\endcsname
	\theoremstyle{plain}
	\newtheorem{theorem}{Theorem}[section]
	\newtheorem{lemma}[theorem]{Lemma}
	\newtheorem{proposition}[theorem]{Proposition}
	\newtheorem*{proposition*}{Proposition}
	%\newtheorem{corollary}[theorem]{אין חלופה עברית}

	\theoremstyle{definition}
	\newtheorem{definition}[theorem]{Definition}
	\newtheorem*{definition*}{Definition}
	\newtheorem{example}{Example}[section]
	\newtheorem{exercise}{Exercise}[section]

	\theoremstyle{remark}
	\newtheorem*{remark}{Remark}
	\newtheorem*{solution}{Solution}
	\newtheorem{conclusion}[theorem]{Conclusion}
	\newtheorem{notation}[theorem]{Notation}
\else
	\theoremstyle{c_plain}
	\newtheorem{theorem}{משפט}[section]
	\newtheorem{lemma}[theorem]{למה}
	\newtheorem{proposition}[theorem]{טענה}
	\newtheorem*{proposition*}{טענה}
	%\newtheorem{corollary}[theorem]{אין חלופה עברית}

	\theoremstyle{c_definition}
	\newtheorem{definition}[theorem]{הגדרה}
	\newtheorem*{definition*}{הגדרה}
	\newtheorem{example}{דוגמה}[section]
	\newtheorem{exercise}{תרגיל}[section]

	\theoremstyle{c_remark}
	\newtheorem*{remark}{הערה}
	\newtheorem*{solution}{פתרון}
	\newtheorem{conclusion}[theorem]{מסקנה}
	\newtheorem{notation}[theorem]{סימון}
\fi

% Questions related commands
\newcounter{question}
\setcounter{question}{1}
\newcounter{sub_question}
\setcounter{sub_question}{1}

\ifcsname c@english\endcsname
	\newcommand{\question}[1][0]{
		\ifthenelse{#1 = 0}{}{\setcounter{question}{#1}}
		\section{Question \arabic{question}}
		\addtocounter{question}{1}
		\setcounter{sub_question}{1}
	}

	\newcommand{\subquestion}[1][0]{
		\ifthenelse{#1 = 0}{}{\setcounter{sub_question}{#1}}
		\subsection{Part \alph{sub_question}}
		\addtocounter{sub_question}{1}
	}
\else
	\newcommand{\question}[1][0]{
		\ifthenelse{#1 = 0}{}{\setcounter{question}{#1}}
		\section{שאלה \arabic{question}}
		\addtocounter{question}{1}
		\setcounter{sub_question}{1}
	}

	\newcommand{\subquestion}[1][0]{
		\ifthenelse{#1 = 0}{}{\setcounter{sub_question}{#1}}
		\subsection{סעיף \localecounter{letters.gershayim}{sub_question}}
		\addtocounter{sub_question}{1}
	}
\fi

% import lua and start of document
\directlua{common = require ('../common')}

\GetEnv{AUTHOR}

% headers
\author{\AUTHOR}
\date\today

\title{פתרון ממ''ן 16 – חשבון אינפיניטסימלי 2 (20475)}

\begin{document}
\maketitle
\maketitleprint{}

\section{שאלה 1}
\subsection{סעיף א'}
נחשב את הגבולות הבאים או נראה כי אינם מתקיימים

\subsubsection{1}
\[
	\lim_{(x, y) \to (0, 0)} (x + y) \ln(x^4 + y^4)
\]
נשים לב כי לכל $|x, y| < 1$ מתקיים $x^4 + y^4 < x + y$.
לכן נובע
\[
	|(x + y) \ln(x^4 + y^4)| \le |x + y| \ln(|x + y|)
\]
נגדיר $t = |x + y|$ ונקבל
\[
	\lim_{(x, y) \to (0, 0)} |x + y| \ln(|x + y|)
	= \lim_{t \to 0^+} t \ln(t)
	= \lim_{t \to 0^+} \frac{\ln(t)}{\frac{1}{t}}
	\overset{\text{לופיטל}}{=} \lim_{t \to 0^+} -x
	= 0
\]
ולכן נקבל מכלל הסנדוויץ' כי גם הגבול המקורי מתכנס לאפס.

\subsubsection{2}
\[
	\lim_{(x, y) \to (0, 0)} \frac{xy^3}{x^2 + y^6}
\]
נניח כי סדרת הנקודות $(p_n = (x_n, y_n))$ מתכנסת ל־$L$,
ולכן מהגדרת היינה לגבולות בשני משתנים נובע שהגבול מתכנס עבור $p_n = (\frac{1}{n}, 0)$:
\[
	\lim_{n \to \infty} f(p_n)
	= \frac{\frac{0^3}{n}}{\frac{1}{n^2} + 0^6}
	= L = 0
\]
כמו־כן, הגבול מתכנס עבור הסדרה $p_n = (\frac{1}{n^3}, \frac{1}{n})$:
\[
	\lim_{n \to \infty} f(p_n)
	= \frac{\frac{1}{n^3 \cdot n^3}}{\frac{1}{n^6} + \frac{1}{n^6}}
	= L = \frac{1}{2}
\]
אז מצאנו כי עבור סדרות שונות הגבול מתכנס לערכים שונים ובהתאם להגדרת היינה הגבול לא מתכנס בנקודה.

\subsection{סעיף ב'}
נבדוק את רציפות הפונקציות הבאות ב־$\RR^2$:

\subsubsection{1}
\[
	f(x, y) = \begin{cases}
		\frac{xy}{x^2 + y^2} & (x, y) \ne (0, 0) \\
		0 & (x, y) = (0, 0) \\
	\end{cases}
\]
לכל נקודה $p_0 = (x_0, y_0)$ כאשר $x_0, y_0 \ne 0$ נראה כי נובע מרציפות והגדרת היינה להתכנסות כי
\[
	\lim_{p \to p_0} f(p)
	= \frac{\lim_{(x, y) \to p_0} xy}{\lim_{(x, y) \to p_0} x^2 + y^2}
	= \frac{x_0y_0}{x_0^2 + y_0^2}
	= f(p_0)
\]
אז מצאנו כי $f$ רציפה לכל $p \ne \textbf{0}$. \\*
מדוגמה 7.13 עולה כי ל־$f$ אין גבול בנקודה $\textbf{0}$ ולכן בהכרח גם איננה מתכנסת בנקודה.

\subsubsection{2}
\[
	\begin{cases}
		\frac{\sin(2x^2 + 2y^2)}{e^{x^2 + y^2} - 1} & (x, y) \ne (0, 0)\\
		2 & (x, y) = (0, 0)
	\end{cases}
\]
נגדיר $t = x^2 + y^2$ ונגדיר פונקציה חדשה $g$ המקיימת:
\[
	g(t) = \begin{cases}
		\frac{\sin(2t)}{e^{t} - 1} & x > 0 \\
		2 & x = 0
	\end{cases}
\]
מתקיים  $g(t) = f(x, y)$ לכל $t$.
$g$ היא כמובן רציפה לכל $x > 0$, ולכן נבדוק את גבולה ב־$0$ בעזרת לופיטל:
\[
	\lim_{x \to \infty} g(x)
	= \lim_{x \to \infty} \frac{\sin(2x)}{e^x - 1}
	= \lim_{x \to \infty} \frac{2\cos(2x)}{e^x}
	= \frac{2}{1}
	= 2
\]
אז מצאנו כי $g$ רציפה לכל $|t| \ge 0$ ובהתאם $f$ רציפה בכל המישור.

\section{שאלה 2}
נמצא את כל הנקודות $(x, y) \in \RR^2$ בהן הפונקציה $f(x, y) = \sqrt{\lvert xy^3 \rvert}$ היא דיפרנציאבילית. \\*
נמצא את הנגזרות החלקיות של $f$ כאשר $xy$ חיובי וכאשר הוא שלילי:
\begin{align*}
	xy \ge 0, f_x(x, y)
	& = \frac{y^3}{2 \sqrt{xy^3}}
	= \frac{y^2}{2 \sqrt{xy}} \\
	xy < 0, f_x(x, y)
	& = \frac{-y^3}{2 \sqrt{-xy^3}}
	= \frac{-y^2}{2 \sqrt{-xy}} \\
	xy \ge 0, f_y(x, y)
	& = \frac{3xy^2}{2 \sqrt{xy^3}}
	= \frac{3xy}{2 \sqrt{xy}}
	= \frac{3}{2} \sqrt{xy} \\
	xy < 0, f_y(x, y)
	& = \frac{-3xy}{2 \sqrt{-xy}}
	= \frac{-3}{2} \sqrt{-xy} \\
\end{align*}
נשים לב כי עבור $x \ne 0, y \ne 0$ כל נקודה שייכת לזוג נגזרות חלקיות ורציפה בהן ובהתאם לתנאי המספיק לדיפרנציאביליות $f$ דיפרנציאבילית בנקודה $(x, y)$. \\*
נבדוק את התכנסות הנגזרות כאשר $y$ הוא ערך קבוע לא אפס ו־$x = 0$: \\*
מונה $f_x$ הוא רציף ושואף לערך סופי, בעוד המכנה שואף לאינסוף לכל $y$ ולכן $f_x$ לא רציפה בנקודות אלה,
וממשפט 7.63 אנו יכולים להסיק כי $f$ לא דיפרנציאבילית בתחום זה. \\*
נבדוק את התכנסות הנגזרות כאשר $x$ הוא ערך קבוע ו־$y = 0$: \\*
בנקודות אלה $f_x$ כמובן רציפה, שכן המונה שלה מתאפס והמכנה רציף ולא מתאפס. כמו־כן גם $f_y$ פונקציה רציפה ומתאפסת באוסף הנקודות.
לכן $f$ דיפרנציאבילית באוסף הנקודות $(x, 0)$. \\*
מצאנו כי $f$ דיפרנציאבילית לכל $\RR^2$ חוץ מהנקודות $(0, t)$ כאשר $t \ne 0$.

\section{שאלה 3}
\subsection{סעיף א'}
נוכיח כי $u(x, y) = \ln(x^2 + y^2)$ היא הרמונית לכל $\RR^2 \backslash \{(0, 0)\}$.
\begin{proof}
	נחשב את נגזרותיה החלקיות מסדר שני של $u$ בתחום:
	\begin{align*}
		u_x(x, y) & = \frac{2x}{x^2 + y^2} \\
		u_{xx}(x, y) & = \frac{2x^2 + 2y^2 - 4x^2}{{(x^2 + y^2)}^2} = 2\frac{y^2 - x^2}{{(x^2 + y^2)}^2} \\
		u_y(x, y) & = \frac{2y}{x^2 + y^2} \\
		u_{yy}(x, y) & = \frac{2x^2 + 2y^2 - 4y^2}{{(x^2 + y^2)}^2} = 2\frac{x^2 - y^2}{{(x^2 + y^2)}^2} \\
	\end{align*}
	נבדוק
	\[
		f_{xx}(x, y) + f_{yy}(x, y)
		= 2\frac{y^2 - x^2}{{(x^2 + y^2)}^2} + 2\frac{x^2 - y^2}{{(x^2 + y^2)}^2}
		= 0
	\]
	ומצאנו כי $f$ הרמונית בתחום $\RR^2\backslash\{(0, 0)\}$.
\end{proof}

\subsection{סעיף ב'}
נגדיר $f(x, y) = u(x^2 - y^2, y^2 - x^2)$ כאשר $u(t, s)$ פונקציה דיפרנציאבילית בכל נקודה. \\*
נוכיח כי $x \, f_y + y \, f_x = 0$ לכל $x, y$.
\begin{proof}
	נגדיר $t(u, v) = u^2 - v^2$, ולכן $f(p) = u(t(p), -t(p))$.
	נחשב את נגזרותיה החלקיות של $f$ על־פי כלל השרשרת המוכלל מהתוספת לחומר של יחידה 7:
	\[
		f_x(x, y)
		= u_x t_u + u_y t_v
		= u_x(x, y) \cdot 2x + u_y(x, y) \cdot (-2y)
		= 2x (u_x(x, y) - u_y(x, y))
	\]
	ובאופן דומה נקבל גם
	\[
		f_y(x, y)
		= u_x t_v + u_y (-t_v)
		= (-2x) u_x - (-2x) u_y
		= 2y (-u_x(x, y) + u_y(x, y))
	\]
	ונקבל כי
	\[
		x f_y + y f_x
		= 2xy (u_x(x, y) - u_y(x, y))
		+ 2yx (-u_x(x, y) + u_y(x, y))
		= 0
	\]
	ומצאנו כי השוויון מתקיים.
\end{proof}

\subsection{סעיף ג'}
נמצא את הקצב בו משתנה שטח מלבן אשר ברגע נתון אורכו 15 מטרים והוא קטן ב־3 מטרים לשנייה, \\*
ואשר רוחבו הוא 6 מטרים והוא גדל ב־2 מטרים לשנייה. \\*
נגדיר $f(x, y) = xy$ פונקציית השטח של מלבן. \\*
עוד נגדיר $h(t)$ פונקציית הגובה של המלבן לפי זמן ו־$w(t)$ פונקציית הרוחב של המלבן לפי זמן. \\*
נגדיר $F(t) = f(h(t), w(t))$ פונקציית השטח של המלבן הנתון לפי זמן. \\*
נגדיר את הרגע הנתון בשאלה בתור $t = 0$ ולכן $h(0) = 15, h'(0) = -3, w(0) = 6, w'(0) = 2$. \\*
נשים לב כי $f_x(x, y) = y, f_y(x, y) = x$ ולכן מכלל השרשרת בשני משתנים נובע כי
\[
	F'(t) = f_x(h(t), w(t)) \cdot h'(t) + f_y(h(t), w(t)) \cdot w'(t)
	= w(t) h'(t) + h(t) w'(t)
\]
נציב
\[
	F'(0) = 6 \cdot (-3) + 15 \cdot 2 = 12
\]

\clearpage
\subsection{סעיף ד'}
תהי $f(x, y)$ דיפרנציאבילית במישור כולו. \\*
נוכיח כי לכל שתי נקודות $p_1 = (x_1, y_1)$ ו־$p_2 = (x_2, y_2)$,
בקטע המחבר בין $p_1$ ל־$p_2$ קיימת נקודה $p_c$ כך שמתקיים
\[
	f(p_2) - f(p_1) = f_x(p_c) (x_2 - x_1) + f_y(p_c) (y_2 - y_1)
\]
\begin{proof}
	תהי $l : \RR \to \RR^2$ העתקה לינארית המוגדרת על־ידי
	\[
		l(t) = p_1 + (p_2 - p_1) t = (x_1 + (x_2 - x_1) t, y_1 + (y_2 - y_1) t)
	\]
	ונגדיר גם $g(t) = f(l(t))$.
	נשים לב כי $l$ היא הישר העובר דרך הנקודות $p_1$ ו־$p_2$, והפונקציה אף מקיימת $l(0) = p_1, l(1) = p_2$. \\*
	בד בבד הפונקציה $g$ היא ''החתך'' של $f$ עבור ערכי הנקודות על הישר $t$, דהינו זו פונקציה ממשתנה אחד. \\*
	ממשפט הערך הממוצע של החשבון הדיפרנציאלי נובע כי קיים ערך $t_c$ אשר מקיים
	\[
		g'(t_c) (1 - 0) = g(1) - g(0) = f(p_2) - f(p_1) \tag{1}
	\]
	כאשר $0 < t_c < 1$. \\*
	נחשב את $g'(t)$ על־פי כלל השרשרת בשני משתנים:
	\begin{align*}
		g'(t)
		= & \frac{d}{dx} f(x_1 + (x_2 - x_1)t, y_1 + (y_2 - y_1)t) \\
		= & f_x(x_1 + (x_2 - x_1)t, y_1 + (y_2 - y_1)t) \cdot (x_1 + (x_2 - x_1) t)' \\
		  & + f_y(x_1 + (x_2 - x_1)t, y_1 + (y_2 - y_1)t) \cdot (y_1 + (y_2 - y_1) t)' \\
		= & f_x(l(t)) (x_2 - x_1) + f_y(l(t)) (y_2 - y_1)
	\end{align*}
	ועל־ידי שימוש בנוסחה זו ב־$(1)$ נקבל
	\[
		f_x(l(t_c)) (x_2 - x_1) + f_y(l(t_c)) (y_2 - y_1) = f(p_2) - f(p_1)
	\]
	נגדיר $p_c = l(t_c)$, ברור כי הנקודה נמצאת בין $p_1$ לבין $p_2$, ונקבל כי
	\[
		f_x(p_c) (x_2 - x_1) + f_y(p_c) (y_2 - y_1) = f(p_2) - f(p_1)
	\]
	והוכחנו את הטענה.
\end{proof}

\section{שאלה 4}
\subsection{סעיף א'}
מטייל עולה על הר שצורתו נתונה על־ידי הנוסחה
\[
	h(x, y) = 1000 - 0.05 x^2 - 0.04 y^2
\]
נתון כי המטייל נמצא בנקודה $(60, 100)$, נמצא את הכיוון עליו ללכת כדי להגיע לפסגה מהר ככל האפשר. \\*
המטייל יגיע לפסגה במהירות הגבוהה ביותר אם ילך בדרך התלולה ביותר, ולכן עלינו למצוא את הכיוון בו הנגזרת הכיוונית היא הגבוהה ביותר. \\*
נשים לב כי
\[
	f_x(x, y) = -0.1x,
	f_y(x, y) = -0.08y
\]
נחשב את הגרדיאנט של $h$:
\[
	\nabla f(60, 100)
	= (f_x(60, 100), f_y(60, 100))
	= (-6, -8)
\]
נחשב את הווקטור הנורמלי:
\[
	u = \frac{(-6, -8)}{\lVert(-6, -8)\rVert} = (-0.6, -0.8)
\]
ולכן הכיוון שעל המטייל ללכת בו הוא $(-0.6, -0.8)$.

\subsection{סעיף ב'}
נמצא את הנקודה על המישור $x + 2y + z = 1$ הקרובה ביותר לראשית הצירים. \\*
נשים לב כי $z = 1 - x - 2y$ ולכן נגדיר פונקציה $z(x, y) = 1 - x - 2y$ ונקבל כי לכל $x, y \in \RR$ הנקודה $(x, y, z(x, y))$ על המישור הנתון. \\*
על־פי נוסחת המרחק בין נקודות (פיתגורס) בין נקודה על המישור לבין נקודת האפס הוא:
\[
	d_0(x, y)
	= \sqrt{x^2 + y^2 + z^2(x, y)}
	= \sqrt{x^2 + y^2 + {(1 - x - 2y)}^2}
\]
כמובן שהפונקציה $d_0$ מקבלת מינימום ומקסימום באותן הנקודות כמו $d_0^2$, ולכן נגדיר $d(x, y) = d_0^2(x, y)$ ונחקור אותה. \\*
נמצא את נקודת המינימום של הפונקציה הנתונה בעזרת משפט 7.58, דהינו נמצא נקודות בהן שתי הנגזרות החלקיות מתאפסות.
\[
	d_x(x, y) = 2x - 2(1 - x - 2y) = -2 + 4x + 4y,
	d_y(x, y) = 2y - 4(1 - x - 2y) = -4 + 4x + 10y,
\]
נמיר למטריצת מקדמים ונפתור
\[
	\begin{pmatrix}
		4 & 4 & \vline & 2 \\
		4 & 10 & \vline & 4
	\end{pmatrix}
	\rightarrow
	\begin{pmatrix}
		4 & 4 & \vline & 2 \\
		0 & 1 & \vline & \frac{1}{3}
	\end{pmatrix}
	\rightarrow
	\begin{pmatrix}
		4 & 0 & \vline & \frac{2}{3} \\
		0 & 1 & \vline & \frac{1}{3}
	\end{pmatrix}
\]
ומצאנו כי יש למערכת פתרון יחיד $(\frac{1}{6}, \frac{1}{3})$. \\*
נחשב את הנגזרת השנייה $d_{xx}(x, y) = 4$, ולכן ממשפט 7.72 נובע כי
$(\frac{1}{6}, \frac{1}{3})$ היא נקודת מינימום של הפונקציה ובה המרחק מהראשית הוא הקטן ביותר על המישור.

\section{שאלה 5}
נוכיח או נפריך את הטענות הבאות:

\subsection{סעיף א'}
נוכיח כי אם $p_0$ נקודת מינימום מקומי של $f(x, y)$ ו־$f$ דיפרנציאבילית בנקודה זו, אז $(D_v f)(p_0) = 0$ עבור כל כיוון $v$.
\begin{proof}
	 ממשפט 7.58 נובע כי $f_x(p_0) = f_y(p_0) = 0$. \\*
	 נגדיר $v = (v_x, v_y)$ וקטור יחידה כלשהו, ממשפט 7.67 נובע כי
	 \[
		 (D_v f)(p_0) = v_x \cdot f_x(p_0) + v_y \cdot f_y(p_0) = 0v_x + 0v_y = 0
	 \]
\end{proof}

\subsection{סעיף ב'}
נסתור את הטענה כי אם $\nabla f(p_0) = (0, 0)$ אז $(D_v f)(p_0) = 0$ עבור כל כיוון $v$.
\begin{proof}[דוגמה נגדית]
	נגדיר $p_0 = 0$ וגם 
	\[
		f(x, y) =
		\begin{cases}
			\frac{xy}{x^2 + y^2} & (x, y) \ne 0 \\
			0 & (x, y) = 0
		\end{cases}
	\]
	בהתבסס על דוגמה 7.13 נגדיר סדרת נקודות $(\frac{1}{n}, 0)$ וניעזר בה לחישוב
	\[
		f_x(0, 0)
		= \lim_{x \to 0} \frac{f(x, 0) - f(0, 0)}{x}
		= \lim_{n \to \infty} \frac{0}{\frac{1}{n^2} + 0^2}
		= 0
	\]
	ובחישוב דומה נקבל גם $f_y(0, 0) = 0$, לכן $\nabla f(p_0) = (0, 0)$. \\*
	נחשב עתה את הנגזרת הכיוונית של $f$ עבור $v = (\sqrt{\frac{1}{2}}, \sqrt{\frac{1}{2}})$ על־ידי הגדרה 7.59 וסדרת הנקודות $(\frac{1}{n}, \frac{1}{n})$:
	\[
		\lim_{t \to 0} \frac{f(p_0 + vt) - f(p_0)}{t}
		= \lim_{n \to \infty} \frac{f(\frac{1}{n}, \frac{1}{n})}{t}
		= \lim_{n \to \infty} \frac{\frac{1}{n^2}}{\frac{1}{n} (\frac{1}{n^2} + \frac{1}{n^2})}
		= \lim_{n \to \infty} \frac{n}{2}
		= \infty
	\]
	ומצאנו כי לא קיימת נגזרת כיוונית ב־$p_0$ לכיוון $v$ בסתירה לתנאי הטענה.
\end{proof}

\subsection{סעיף ג'}
נפריך את הטענה כי אם פונקציה $f(x, y)$ רציפה בקבוצה $A = \{ (x, y) : x^2 + y^2 \ge 1 \}$ אז יש ל־$f$ מינימום ב־$A$.
\begin{proof}[דוגמה נגדית]
	נגדיר $f(p) = \frac{1}{\lVert p \rVert}$. \\*
	פונקציה זו כמובן מוגדרת ב־$\RR^2\backslash \{0\}$ ובפרט ב־$A$. \\*
	נשים לב כי נקודה $p = (x, y)$ בקבוצה $A$ אם ורק אם $\lVert (x, y) \rVert \ge 1$,
	ונראה כי גם $p_1 = (x + 1, y + 1)$ מקיימת $\lVert p_1 \rVert > \lVert p \rVert \ge 1$ ולכן מהגדרת $f$ גם $f(p_1) < f(p)$, \\*
	ומצאנו לכל נקודה של $f$ ב־$A$ נקודה עם ערך קטן ממנה, דהינו אין ל־$f$ מינימום ב־$A$ כלל בסתירה לטענה.
\end{proof}

\end{document}
